%% Drafted from the mapped corpus by scripts/draft_topics.py.
% Model: gpt-5.6-luna; source characters: 10195.
\begin{konspekt}
\kreq{Официално заглавие}{„Компютърни архитектури --- Формати на данните. Вътрешна структура на централен процесор --- \emph{блокове и конвейерна обработка, инструкции}.“ Заглавието на главата е съкратено; конвейерът е изискване, не допълнение.}
\kreq{Описание}{обща структура на компютрите, запомнена програма --- \kref{9.1}.}
\kreq{Описание}{концептуално изпълнение на инструкциите --- \kref{9.4}.}
\kreq{Формати на данните}{цели двоични числа, двоично-десетични числа, числа с плаваща запетая, знакови данни и кодови таблици --- \kref{9.2}.}
\kreq{Централен процесор}{регистри, АЛУ, регистри на състоянията и флаговете, блокове за управление, връзка с паметта, дешифрация на инструкциите, преходи --- \kref{9.3}.}
\kreq{Конвейерна обработка}{\kref{9.5}.}
\kreq{Литература}{[5], [36], [45].}
\end{konspekt}

\section{Обща структура на компютрите и компютърни архитектури}\label{s:9.1}

Компютърната архитектура описва основните компоненти на компютъра, начина на свързването им и принципите, по които се изпълняват програмите. В общия случай компютърът се състои от:

\begin{itemize}
    \item централен процесор;
    \item основна памет;
    \item входно-изходни устройства.
\end{itemize}

Тези компоненти са взаимносвързани и работят съвместно, за да изпълняват програми. Централният процесор обработва данните и управлява изпълнението на инструкциите. Основната памет съхранява инструкции и данни. Входно-изходните устройства осигуряват обмен на информация между компютъра и външната среда.

Връзката между компонентите се осъществява чрез шини.

\begin{defn}[Шина]
Шината е съвкупност от проводници, по които се предава информация между компонентите на компютъра. Обичайно се разграничават три шини:

\begin{itemize}
\item \emph{адресна шина} --- по нея процесорът подава адреса на клетката от паметта или на устройството, към което се обръща;
\item \emph{шина за данни} --- по нея се пренася самата стойност, която се чете или записва;
\item \emph{контролна шина} --- по нея се предават управляващите сигнали, например дали текущото обръщение е четене или запис, и сигналите за прекъсване.
\end{itemize}
\end{defn}

\subsection{Архитектура на фон Нойман}\label{s:9.1.1}

Повечето съвременни компютри се основават на архитектурата на Джон фон Нойман. Тя се характеризира със следните основни идеи:

\begin{enumerate}
    \item Инструкциите и данните се съхраняват в памет, която позволява четене и запис.
    \item Съдържанието на паметта се адресира чрез адрес, независимо от вида на данните, които се намират на съответното място.
    \item Инструкциите се изпълняват последователно, от една инструкция към следващата, освен ако текущата инструкция не промени реда на изпълнение.
\end{enumerate}

Последната особеност е съществена за разбирането на програмния брояч. Той съдържа адреса на инструкцията, която трябва да бъде извлечена и изпълнена. При обикновено последователно изпълнение програмният брояч се насочва към следващата инструкция. При преход или друга инструкция, която променя управлението, в него се записва различен адрес.

\begin{defn}[Запомнена програма]
Запомнена програма е програма, чиито инструкции са представени в двоичен вид и са съхранени в паметта заедно с данните, върху които работи програмата. Процесорът извлича инструкциите от паметта и ги изпълнява последователно, като редът може да бъде променен чрез инструкции за преход.
\end{defn}

Така програмата не е външна поредица от действия, която процесорът следва по специален начин, а е информация, записана в паметта. Процесорът интерпретира тази информация като инструкции и данни според съответния формат.

\section{Формати на данните}\label{s:9.2}

\begin{defn}[Формат на данни]
Форматът на данните определя каква битова последователност представя дадена стойност. Една и съща последователност от битове може да има различно значение в зависимост от това дали се разглежда като цяло число, двоично-десетично число, число с плаваща запетая или код на символ.
\end{defn}

\subsection{Цели двоични числа}\label{s:9.2.1}

Целите двоични числа се представят в двоична бройна система. При знаковите числа един от битовете обикновено се използва за представяне на знака. Разглеждат се различни кодове за представяне на отрицателните числа.

\subsubsection{Прав код}

\begin{defn}[Прав код]
При правия код най-старшият бит в $n$-битовото поле е знаков. Стойност $0$ означава положително число, а стойност $1$ означава отрицателно число. Останалите битове представят абсолютната стойност на числото в двоична бройна система.
\end{defn}

Следователно в прав код има отделно поле за знака и отделно поле за абсолютната стойност. Основните недостатъци на този начин на представяне са:

\begin{itemize}
    \item нулата има две представяния -- положителна и отрицателна;
    \item събирането и изваждането изискват по-сложна апаратна реализация.
\end{itemize}

Диапазонът на представимите стойности е от $-(2^{n-1}-1)$ до $2^{n-1}-1$.

\subsubsection{Двоично-допълнителен код}

\begin{defn}[Двоично-допълнителен код]
При двоично-допълнителния код $n$-битовата последователност $b_{n-1}b_{n-2}\ldots b_0$ представя числото
\[
-b_{n-1}2^{n-1}+\sum_{i=0}^{n-2}b_i2^i .
\]
Тоест всички битове имат обичайните си тегла, но най-старшият участва с отрицателно тегло. Затова той отново определя знака: стойност $1$ дава отрицателно число.

Еквивалентно, отрицателното число $x<0$ се записва като двоичния запис на $2^n+x$; оттук идва и името на кода.
\end{defn}

От определението следва диапазонът на представимите стойности:
\[
[-2^{n-1},\,2^{n-1}-1],
\]
защото най-малката стойност се получава при $b_{n-1}=1$ и всички останали битове нула, а най-голямата --- при $b_{n-1}=0$ и всички останали единици.

Този код използва едно представяне на нулата и позволява операциите със знакови числа да се реализират чрез обикновено двоично събиране, като се работи с фиксиран брой битове.

Смяната на знака на число в двоично-допълнителен код се извършва в две стъпки:

\begin{enumerate}
    \item всички битове се инвертират побитово;
    \item към получения резултат се прибавя $1$.
\end{enumerate}

\begin{prop}
Нека $x$ е представено в $n$-битов двоично-допълнителен код и нека $\overline{x}$ е побитовото му инвертиране. Тогава $\overline{x}+1$ е представянето на $-x$ при пресмятане с $n$ бита.
\end{prop}

\begin{proof}
Инвертирането на бит го заменя с $1$ минус стойността му, затова при разглеждане на битовете като безнакови тегла
\[
x+\overline{x}=\underbrace{11\ldots1}_{n}=2^n-1 .
\]
Следователно
\[
\overline{x}+1=2^n-x .
\]
При работа с $n$ бита преносът извън най-старшия бит се отхвърля, тоест пресмята се по модул $2^n$, а
\[
2^n-x\equiv -x \pmod{2^n}.
\]
\end{proof}

\begin{example}
Нека числото $3_{10}$ бъде представено с четири бита:

\[
3_{10}=0011_2.
\]

Инвертирането на битовете дава

\[
0011 \longrightarrow 1100.
\]

След прибавяне на $1$ се получава

\[
1100+1=1101.
\]

Следователно $1101_2$ е представянето на $-3_{10}$ в четирибитов двоично-допълнителен код.
\end{example}

При събиране на число с противоположното му число резултатът е нула в рамките на използвания брой битове. Преносът извън най-старшия бит не се включва в резултата.

\begin{supp}[Бележка]
При работа с двоично-допълнителен код броят на битовете трябва да бъде фиксиран предварително. Например в четирибитово поле резултатът от събиране може да съдържа пренос извън полето, който не е част от представената стойност.
\end{supp}

\subsection{Двоично-десетични числа}\label{s:9.2.2}

\begin{defn}[Двоично-десетичен код]
При двоично-десетичното представяне всяка десетична цифра се заменя с уникална двоична последователност. За една десетична цифра се използват четири бита. Понеже с четири бита могат да се образуват $16$ комбинации, а десетичните цифри са само десет, шест комбинации остават неизползвани.
\end{defn}

Предимството на двоично-десетичното представяне е, че числата са удобни за извеждане и за работа с десетични цифри. Недостатък е усложненото изпълнение на аритметичните операции.

Различават се две основни форми:

\begin{description}
    \item[Непакетирана форма.] Всяка десетична цифра се съхранява в отделен байт.
    \item[Пакетирана форма.] В един байт се съхраняват две десетични цифри, като всяка цифра заема по четири бита.
\end{description}

При непакетираната форма се използва повече памет, но обработката на отделните цифри е по-пряка. При пакетираното представяне паметта се използва по-икономично, защото два четирибитови кода се съхраняват в един байт.

\subsection{Двоични числа с плаваща запетая}\label{s:9.2.3}

Числата с плаваща запетая се използват за представяне на нецели числа. Общата идея е стойността да се представи чрез знак, мантиса и характеристика, тоест като произведение на мантиса и степен на основата:

\[
\pm S \times B^{\pm E}.
\]

\begin{defn}[Мантиса, характеристика, отклонение]
В представянето $\pm S\times B^{\pm E}$:

\begin{itemize}
\item $S$ е \emph{мантисата} --- значещите цифри на числото;
\item $B$ е \emph{основата} на представянето, при двоичните формати $B=2$;
\item $E$ е \emph{характеристиката}, наричана още \emph{експонента} --- степента, на която се повдига основата, тоест с колко позиции се измества запетаята.
\end{itemize}

Полето за характеристиката съдържа не самата стойност $E$, а изместена стойност, за да се представят и отрицателни експоненти без отделен знаков бит. Фиксираната добавка се нарича \emph{отклонение}: истинската експонента се получава, като отклонението се извади от записаното в полето.
\end{defn}

В един двоичен формат числото може да бъде разделено на три части:

\begin{itemize}
    \item един бит за знака;
    \item поле за характеристиката, или експонентата;
    \item поле за мантисата.
\end{itemize}

Посоченият пример използва формат с общо $32$ бита:

\begin{itemize}
    \item $1$ бит за знака;
    \item $8$ бита за характеристиката;
    \item $23$ бита за мантисата.
\end{itemize}

За характеристиката обикновено се използва отклонение. Отклонението е фиксирана стойност, която се изважда от записаната характеристика, за да се получи действителната стойност на експонентата. При $k$ бита за характеристиката обичайно се използва отклонение

\[
2^{k-1}-1.
\]

При $k=8$ това отклонение е $127$. Източникът посочва диапазон на действителните стойности на характеристиката от $-127$ до $128$.

За да се избегнат различни представяния на една и съща стойност, числата се нормализират. Без такова условие едно и също число има много записи: например
\[
1{,}5=1{,}5\times2^0=0{,}75\times2^1=0{,}375\times2^2=\ldots
\]

\begin{defn}[Нормализирано число]
Ненулево число с плаваща запетая е \emph{нормализирано}, ако характеристиката му е избрана така, че мантисата да попада в предварително фиксиран интервал. При основа $B=2$ това означава
\[
1\leq S<2,
\]
тоест числото се записва във вида
\[
\pm 1{,}f\times 2^{E}.
\]
При това условие всяко ненулево число има точно едно представяне.
\end{defn}

Тъй като при двоична основа старшата цифра на нормализирана мантиса е винаги $1$, тя не се записва в полето за мантисата; говори се за \emph{скрит бит}. Така при $23$ бита се получава точност от $24$ значещи двоични цифри. Нулата не може да се нормализира и за нея се запазва отделно представяне.

\begin{supp}[Бележка]
Изречението, което определя нормализираното число, е останало незавършено в източника: формулата след него липсва. Дадената по-горе форма е стандартната за двоичните формати и е добавена, за да бъде определението използваемо.
\end{supp}

Аритметиката на числата с плаваща запетая е по-сложна от аритметиката на целите числа. Стандартното аритметико-логическо устройство не е достатъчно за непосредственото изпълнение на всички операции върху тях. Затова се използва отделно изчислително звено -- блок за плаваща запетая, наричан още \emph{Floating Point Unit}.

Инструкциите за работа с числа с плаваща запетая имат специфични особености. Поради това те се разглеждат отделно от инструкциите за обикновени цели двоични числа.

\begin{supp}[Бележка]
При числата с плаваща запетая точният диапазон и конкретните специални представяния зависят от избрания формат. Тук се разглежда само общата структура, дадена в източника: знак, характеристика и мантиса.
\end{supp}

\subsection{Знакови данни и кодови таблици}\label{s:9.2.4}

\begin{defn}[Кодова таблица]
Кодова таблица е съответствие между символите от определена азбука и битови последователности с определена дължина, чрез които тези символи се съхраняват и обработват в компютъра.
\end{defn}

Кодовите таблици позволяват на процесора и на програмите да обработват букви, цифри, препинателни знаци и други символи като цифрови данни.

Източникът посочва два стандарта:

\begin{description}
    \item[ASCII.] Използва седем бита за кодиране на английската азбука, цифрите, пунктуацията и някои други символи.
    \item[Unicode.] Стандарт за представяне на символи чрез кодови стойности, предназначен за работа с различни азбуки и знакови системи.
\end{description}

\section{Централен процесор}\label{s:9.3}

Централният процесор изпълнява инструкциите на програмата и управлява обработката на данните. За тази цел той трябва да извършва няколко основни действия:

\begin{enumerate}
    \item да извлича инструкции от паметта;
    \item да интерпретира инструкциите;
    \item да извлича необходимите операнди;
    \item да обработва данните;
    \item да записва резултатите.
\end{enumerate}

Извличането на инструкция означава прочитане на инструкция от достъпно за процесора място -- регистри, кеш или основна памет. Интерпретирането включва дешифрация на инструкцията, за да се установи каква операция трябва да бъде изпълнена и какви са нейните операнди.

Изпълнението може да изисква четене на данни от паметта или от входно-изходен модул. Самата обработка може да представлява аритметична или логическа операция. След приключването й резултатът се записва в регистър, памет или входно-изходно устройство.

Основните компоненти на процесора са:

\begin{itemize}
    \item аритметико-логическо устройство;
    \item блок за управление;
    \item регистри.
\end{itemize}

\subsection{Регистри}\label{s:9.3.1}

\begin{defn}[Регистър]
Регистрите са вътрешни места за временно съхраняване на данни, адреси, инструкции и информация за състоянието на процесора. Те позволяват на процесора да работи с междинни резултати и да запомня позицията си в програмата.
\end{defn}

Регистрите могат да се разделят на две основни групи:

\begin{enumerate}
    \item потребителски видими регистри;
    \item регистри за управление и състояние.
\end{enumerate}

\subsubsection{Потребителски видими регистри}

\begin{defn}[Потребителски видим регистър]
Потребителски видими са регистрите, които могат да бъдат достъпвани чрез машинни инструкции. Те позволяват на програмиста или на компилатора да съхранява междинни резултати и да намалява необходимостта от обръщения към основната памет.
\end{defn}

Основните подкатегории са:

\begin{description}
    \item[Регистри с общо предназначение.] Съхраняват междинни резултати и могат да се използват при адресиране.
    \item[Регистри за данни.] Предназначени са за съхраняване на данни, но не и за изчисляване на адрес на операнд.
    \item[Регистри за адреси.] Използват се за адресиране. Те могат да бъдат регистри с общо предназначение или да са предназначени за конкретен режим на адресиране, например като индексни регистри, указатели на стека или сегментни указатели.
\end{description}

Границата между потребителски видимите регистри и регистрите за управление и състояние не е еднаква при всички процесори. Някои регистри, използвани основно от операционната система, могат да бъдат достъпни чрез машинни инструкции, макар да не се използват пряко от обикновените програми.

\subsubsection{Регистри за управление и състояние}

Четири регистъра са особено важни за изпълнението на инструкциите:

\begin{description}
    \item[Програмен брояч.] Съдържа адреса на инструкцията, която трябва да бъде прочетена. При последователно изпълнение този адрес се променя така, че да сочи към следващата инструкция.
    \item[Регистър на инструкциите.] Съдържа последно извлечената инструкция, която се декодира и изпълнява.
    \item[Регистър на адреса на паметта (MAR).] Съдържа адреса на място в паметта, към което ще се извърши обръщение.
    \item[Буферен регистър на паметта (MBR).] Съдържа данните, които трябва да бъдат записани в паметта, или последно прочетените от паметта данни.
\end{description}

Много процесори имат и регистър на флаговете.

\begin{defn}[Флаг и регистър на състоянието]
\emph{Флаг} е отделен бит, който записва едно двоично условие --- най-често свойство на резултата от последната операция на аритметико-логическото устройство или режим на работа на процесора. Регистърът, в който са събрани тези битове, се нарича \emph{регистър на състоянието} или \emph{регистър на флаговете}.

Флаговете не се четат като обикновени данни, а се използват от инструкциите за условен преход: те решават коя да бъде следващата инструкция.
\end{defn}

Сред често срещаните флагове са:

\begin{itemize}
    \item флаг за знак;
    \item флаг за нулев резултат;
    \item флаг за пренос;
    \item флаг за препълване;
    \item флаг за разрешени или забранени прекъсвания.
\end{itemize}

Флагът за знак показва знака на резултата, флагът за нула показва дали резултатът е нулев, а флаговете за пренос и препълване описват определени условия при аритметични операции. Флагът за прекъсване участва в управлението на възможността процесорът да приема прекъсвания.

\subsection{Аритметико-логическо устройство}\label{s:9.3.2}

\begin{defn}[Аритметико-логическо устройство]
Аритметико-логическото устройство, съкратено АЛУ, извършва аритметични и логически операции върху операнди. То получава данни от регистрите или от пътя за данни, изпълнява избраната операция и връща резултата към регистър или към следващ етап на обработката.
\end{defn}

Работата на АЛУ се определя от управляващи сигнали. Те задават каква операция трябва да бъде извършена и как резултатът да бъде прехвърлен. В зависимост от резултата АЛУ може да промени и флаговете в регистъра на състоянието.

\subsection{Блок за управление}\label{s:9.3.3}

\begin{defn}[Микрооперация]
Микрооперацията е най-малкото неделимо действие, което процесорът извършва за един такт --- например прехвърляне на съдържанието на един регистър в друг, подаване на адрес към паметта или активиране на конкретна функция на аритметико-логическото устройство. Изпълнението на една машинна инструкция се разлага на последователност от микрооперации.
\end{defn}

Блокът за управление координира работата на процесора. Той изпълнява две основни задачи:

\begin{enumerate}
    \item задвижва процесора през поредица от микрооперации в правилната последователност;
    \item генерира управляващи сигнали, които предизвикват изпълнението на микрооперациите.
\end{enumerate}

Управляващите сигнали контролират прехвърлянето на данни между регистрите, избора на операция на АЛУ и обмена с паметта и входно-изходните устройства. Те управляват отварянето и затварянето на логически пътища, чрез които данните достигат до съответните регистри и функционални блокове.

Входове за блока за управление са:

\begin{itemize}
    \item тактовият сигнал;
    \item съдържанието на регистъра на инструкциите;
    \item флаговете на процесора;
    \item управляващи сигнали от контролната шина.
\end{itemize}

Тактовият сигнал задава последователността на микрооперациите. Регистърът на инструкциите предоставя кода на операцията и информацията за режима на адресиране. Флаговете показват състоянието на процесора и резултатите от предишни операции. Управляващите сигнали от контролната шина осигуряват информация за състоянието на паметта и входно-изходните устройства.

Изходите на блока за управление са:

\begin{itemize}
    \item управляващи сигнали вътре в процесора;
    \item управляващи сигнали към контролната шина.
\end{itemize}

Сигналите вътре в процесора осигуряват прехвърляне на данни от един регистър в друг и активират конкретни функции на АЛУ. Сигналите към контролната шина управляват обмена с паметта и входно-изходните устройства.

\subsection{Инструкции}\label{s:9.3.4}

Машинната инструкция е двоично кодирана команда, която определя операцията и необходимите за нея данни. Обикновено една инструкция може да съдържа:

\begin{itemize}
    \item код на операцията;
    \item указател към един или повече входни операнди;
    \item указател към изходен операнд;
    \item указател към следващата инструкция.
\end{itemize}

\begin{defn}[Код на операцията]
Кодът на операцията, или \emph{opcode}, е двоична част от инструкцията, която определя какво действие трябва да извърши процесорът.
\end{defn}

\begin{defn}[Режим на адресиране]
Режимът на адресиране определя как от полето на инструкцията се получава действителният адрес, или самата стойност, на операнда. Сред основните режими са:

\begin{itemize}
\item \emph{непосредствен} --- операндът е записан в самата инструкция;
\item \emph{регистров} --- операндът е в посочения регистър;
\item \emph{директен} --- инструкцията съдържа адреса на операнда в паметта;
\item \emph{недиректен} --- инструкцията съдържа адрес, на който е записан адресът на операнда, така че е необходимо още едно обръщение към паметта.
\end{itemize}
\end{defn}

Операндите могат да се намират в регистри, в паметта или във входно-изходно устройство. За да ги използва, процесорът трябва да определи местоположението им според зададения режим на адресиране.

\subsection{Връзка с паметта}\label{s:9.3.5}

Връзката между процесора и паметта се реализира посредством шини, сред които особено значение има контролната шина. Процесорът подава управляващи сигнали, които указват дали се извършва четене или запис, а чрез регистрите MAR и MBR се задават адресът и обменяните данни.

При четене адресът на необходимото място се поставя в MAR. След осъществяване на обръщението прочетените данни се приемат в MBR. При запис адресът се поставя в MAR, а данните за запис -- в MBR. Блокът за управление подава съответния управляващ сигнал към паметта.

\section{Концептуално изпълнение на инструкциите}\label{s:9.4}

Състоянието на процесора образува контекста, необходим за изпълнението на следващата инструкция. Важна част от този контекст е програмният брояч, който съдържа адреса на инструкцията, която трябва да бъде извлечена.

На концептуално ниво изпълнението на инструкцията включва следните фази:

\begin{enumerate}
    \item извличане;
    \item декодиране и интерпретиране;
    \item изчисляване на адресите на операндите;
    \item извличане на операндите;
    \item изпълнение;
    \item записване на резултата;
    \item проверка и обслужване на прекъсване.
\end{enumerate}

\subsection{Извличане на инструкция}\label{s:9.4.1}

При извличането процесорът използва адреса от програмния брояч. Този адрес се подава към паметта, а инструкцията се прочита и се поставя в регистъра на инструкциите. След извличането програмният брояч обикновено се насочва към следващата инструкция.

\subsection{Декодиране}\label{s:9.4.2}

Блокът за управление анализира съдържанието на регистъра на инструкциите. От кода на операцията се определя какво действие трябва да бъде извършено. Разпознават се също операндите и режимът на адресиране.

Дешифрацията на инструкциите е необходима, защото една и съща обща процедура трябва да обслужва различни операции -- например аритметични операции, логически операции, прехвърляния на данни и преходи.

\subsection{Изчисляване и извличане на операнди}\label{s:9.4.3}

Ако инструкцията използва операнд от паметта или от входно-изходно устройство, процесорът изчислява адреса на операнда. След това извършва необходимото обръщение и извлича данните.

Ако операндът вече се намира в регистър, допълнително обръщение към основната памет не е необходимо. Във всички случаи блокът за управление трябва да осигури подаването на правилните данни към АЛУ или към друг функционален блок.

\subsection{Изпълнение и записване на резултата}\label{s:9.4.4}

След извличането на операндите се изпълнява операцията, зададена от инструкцията. При аритметична или логическа операция АЛУ обработва данните. Резултатът може да бъде записан в регистър, в паметта или към входно-изходно устройство.

Ако операцията променя състоянието на процесора, се променят и съответните флагове. След това процесорът продължава със следващата инструкция или изпълнява преход, ако инструкцията е променила програмния брояч.

\subsection{Недиректна стъпка}\label{s:9.4.5}

След извличането на инструкцията може да бъде необходимо да се извърши недиректна стъпка. Тя се прилага, когато инструкцията използва недиректна адресация. В такъв случай извлечената от инструкцията информация не е непосредствено крайният адрес на операнда. Процесорът трябва да извърши допълнително обръщение, за да получи необходимия адрес.

\subsection{Прекъсвания}\label{s:9.4.6}

\begin{defn}[Прекъсване]
Прекъсването е сигнал, който кара процесора да спре временно изпълнението на текущата програма и да премине към предварително определена обслужваща процедура. Сигналът може да идва от входно-изходно устройство, от таймер или от самото изпълнение на инструкция, например при опит за деление на нула.

Прекъсването не се проверява по средата на инструкция: то се обслужва след завършването на текущата инструкция, което е и последната фаза от концептуалното изпълнение.
\end{defn}

Ако прекъсванията са разрешени и възникне прекъсване, процесорът запазва текущото състояние на програмата и преминава към обслужване на прекъсването.

Запазеното състояние включва:

\begin{itemize}
    \item съдържанието на програмния брояч, тоест адреса на следващата инструкция;
    \item релевантна информация за изпълняваната програма;
    \item състоянието на регистрите;
    \item състоянието на входно-изходните устройства.
\end{itemize}

След обслужването на прекъсването запазеното състояние позволява изпълнението на прекъснатата програма да бъде продължено.

\subsection{Преходи}\label{s:9.4.7}

При обикновено изпълнение програмният брояч сочи към следващата инструкция. Инструкцията за преход променя това поведение, като записва нов адрес в програмния брояч. Новият адрес може да зависи от флаговете на процесора.

Например при инструкция от вида \verb|increment-and-skip-if-zero| блокът за управление увеличава програмния брояч, ако флагът за нулев резултат е установен. Така флаговете могат да участват пряко в избора на следващата инструкция.

\begin{supp}[Бележка]
Формулировката е тази от източника. Става дума за \emph{допълнително} увеличаване: програмният брояч се увеличава при всяко извличане на инструкция, а при изпълнена проверка се увеличава още веднъж, с което следващата инструкция се прескача. Оттам идва и \emph{skip} в името на инструкцията.
\end{supp}

\section{Конвейерна обработка}\label{s:9.5}

Изпълнението на една инструкция може да бъде представено като последователност от фази: извличане, декодиране, извличане на операнди, изпълнение и записване на резултата. В процесора тези фази се реализират чрез поредица от микрооперации, управлявани от блока за управление.

\begin{defn}[Конвейерна обработка]
При конвейерната обработка изпълнението на инструкциите се разделя на последователни етапи, всеки от които се обслужва от отделна част на процесора. Докато една инструкция преминава през даден етап, следващата вече може да се обработва в предходния. По този начин отделните части на процесора работят едновременно върху различни инструкции.

Конвейерът не съкращава времето за изпълнение на една инструкция, а увеличава броя инструкции, завършвани за единица време.
\end{defn}

Основата на конвейерната обработка е съгласуваното управление на етапите чрез тактовия сигнал. Блокът за управление генерира управляващите сигнали в правилната последователност, така че данните да преминават между регистрите, АЛУ, паметта и останалите блокове.

При инструкции за преход нормалната последователност може да бъде променена, тъй като програмният брояч получава нов адрес. Затова преходите зависят от състоянието на процесора и от резултатите на предходни операции.

\section{Изпитен фокус}\label{s:9.6}

При подготовка за устен изпит трябва да могат да бъдат възпроизведени:

\begin{itemize}
    \item трите основни компонента на компютъра -- централен процесор, основна памет и входно-изходни устройства;
    \item трите основни идеи на архитектурата на фон Нойман;
    \item понятието запомнена програма;
    \item представянето на цели числа в прав и двоично-допълнителен код;
    \item алгоритъмът за смяна на знака в двоично-допълнителен код и защо той работи: $x+\overline{x}=2^n-1$, откъдето $\overline{x}+1=2^n-x\equiv-x \pmod{2^n}$;
    \item пакетираното и непакетираното двоично-десетично представяне;
    \item частите на числото с плаваща запетая -- знак, характеристика и мантиса --- и ролята на отклонението при характеристиката;

    \item какво означава нормализирано число и защо старшият бит на мантисата може да не се записва;
    \item ролята на кодовите таблици, ASCII и Unicode;
    \item основните блокове на процесора -- регистри, АЛУ и блок за управление;
    \item предназначението на програмния брояч, регистъра на инструкциите, MAR и MBR;
    \item флаговете за знак, нула, пренос, препълване и прекъсване;
    \item задачите и входовете и изходите на блока за управление;
    \item съставът на машинната инструкция -- opcode и указатели към операнди и следваща инструкция;
    \item трите шини --- адресна, за данни и контролна --- и ролята на контролната шина при връзката с паметта;

    \item понятието режим на адресиране и разликата между директно и недиректно адресиране;

    \item понятието микрооперация и връзката му с работата на блока за управление;
    \item фазите на изпълнение на инструкцията;
    \item изчисляването на адреси и извличането на операнди;
    \item какво е прекъсване и защо се обслужва след завършването на текущата инструкция;

    \item запазването на състоянието при прекъсване;
    \item ролята на програмния брояч и флаговете при изпълнение на преходи;
    \item идеята на конвейерната обработка и последователността от микрооперации.
\end{itemize}
